% Highlights for the Elsevier IP&M submission of the Sieve manuscript (v2).
% 5 bullets, each at most 85 characters including spaces (rendered text).
% Verified against paper-journal-v2 numbers.tex (2026-10-01, e870b0d):
%   - budgeted question-conditioned information selection formulation
%   - model-free, training-free, 50.9-54.3 ms at 6k words on one CPU core
%   - 3 datasets, 1600 recall records + end-to-end QA n=100 x 7 methods
%   - BM25 end-to-end retain 76.0% vs Sieve 61.8%; Sieve matched with head
%   - reader-side positional profile re-measured (flat at ~1k words)
\documentclass[12pt]{article}
\usepackage[T1]{fontenc}
\usepackage[margin=1in]{geometry}
\usepackage{microtype}
\usepackage{enumitem}
\pagestyle{empty}

\newcommand{\PaperTitle}{Sieve: Budgeted Question-Conditioned
Information Selection for Training-Free Long-Context Compression}

\begin{document}

\begin{center}
{\LARGE\bfseries Highlights}\\[0.6em]
{\large \PaperTitle{}\\[0.35em]}
{\normalsize\itshape Manuscript submitted to Information Processing \& Management
(Full Length Article)}
\end{center}

\vspace{1.2em}

\begin{itemize}[leftmargin=1.6em, itemsep=0.55em]
\item Long-context compression framed as budgeted question-conditioned selection.
\item A training-free, model-free selector runs in milliseconds on one CPU core.
\item Three-dataset frontier study spans accuracy, compression, and compressor cost.
\item End-to-end QA with a production reader: seven methods, 100 records each.
\item Reader-side positional profile is re-measured, not borrowed from prior work.
\end{itemize}

\end{document}
